% GENERATED FILE. Edit canonical lessons, then run npm run generate:books.
% canonical-source-hash: fnv1a64:2fcea607071a934d
% canonical-lessons: FR-C295-avouer, FR-C295-aveu, FR-C295-complice, FR-C295-acquitter, FR-C295-relacher

\chapter{Les aveux}
\label{ch:fr-les-aveux}
\hlchaptermodality{295}

\begin{chapteropening}
\textbf{By the end of this chapter:} \emph{I can say that someone confessed, had an accomplice, was acquitted or was let go.}

Complete the last lesson of chapter 295: I can say that someone confessed, had an accomplice, was acquitted or was let go.
\end{chapteropening}

\section[avouer]{avouer — to confess, to admit}

\label{lesson:FR-C295-avouer}



\begin{quote}
Before the new one: say the French for police custody, then the French for to question.
\end{quote}

\subsection*{You'll want to know: avouer}
\textbf{avouer} — \textquotedblleft{}to confess, to admit\textquotedblright{}.

\begin{quote}
Il a fini par avouer.
\end{quote}

\begin{quote}
\emph{In the end he confessed.}
\end{quote}

\emph{Je dois avouer que…} — I must admit that…

\begin{cousinweb}
\emph{avouer} is inherited through Old French \emph{avoer} from Latin \emph{advocāre} \textquotedblleft{}to call to, to call upon\textquotedblright{} (EtymDB).
\end{cousinweb}

\begin{grammarlens}[title={What you've built}]
One more word for telling, explaining and describing at B1.
\end{grammarlens}

\subsection*{Your turn}
\begin{itemize}
\raggedright
  \item \emph{Say it:} \emph{avouer}
  \item \emph{Say it:} \emph{avouer}, once more
  \item \emph{Say it:} a sentence of your own with \emph{avouer}
  \item \emph{Recall:} say the French for police custody, then the French for to question, then say \emph{avouer} again
\end{itemize}

\begin{tcolorbox}[breakable,colback=teal!4,colframe=teal!35!black,title={Before you move on}]
What does avouer mean? (\textquotedblleft{}To confess, to admit\textquotedblright{}.) Say it once more.
\end{tcolorbox}
\section[l'aveu]{l'aveu — the confession, the avowal}

\label{lesson:FR-C295-aveu}



\begin{quote}
Before the new one: say the French for to question, then the French for to confess.
\end{quote}

\subsection*{You'll want to know: l'aveu}
\textbf{l'aveu} — \textquotedblleft{}the confession, the avowal\textquotedblright{}.

\begin{quote}
Ses aveux ont surpris tout le monde.
\end{quote}

\begin{quote}
\emph{His confession surprised everyone.}
\end{quote}

Often plural: \emph{passer aux aveux} — to confess.

\begin{cousinweb}
\emph{aveu} is built on \emph{avouer} \textquotedblleft{}to confess\textquotedblright{}, which is inherited from Latin \emph{advocāre} \textquotedblleft{}to call to\textquotedblright{} (Lewis \& Short).
\end{cousinweb}

\begin{grammarlens}[title={What you've built}]
One more word for telling, explaining and describing at B1.
\end{grammarlens}

\subsection*{Your turn}
\begin{itemize}
\raggedright
  \item \emph{Say it:} \emph{l'aveu}
  \item \emph{Say it:} \emph{l'aveu}, once more
  \item \emph{Say it:} a sentence of your own with \emph{l'aveu}
  \item \emph{Recall:} say the French for to question, then the French for to confess, then say \emph{l'aveu} again
\end{itemize}

\begin{tcolorbox}[breakable,colback=teal!4,colframe=teal!35!black,title={Before you move on}]
What does l'aveu mean? (\textquotedblleft{}The confession, the avowal\textquotedblright{}.) Say it once more.
\end{tcolorbox}
\section[le complice]{le complice — the accomplice}

\label{lesson:FR-C295-complice}



\begin{quote}
Before the new one: say the French for to confess, then the French for the confession.
\end{quote}

\subsection*{You'll want to know: le complice}
\textbf{le complice} — \textquotedblleft{}the accomplice\textquotedblright{}.

\begin{quote}
Il avait un complice qui l'attendait dans la voiture.
\end{quote}

\begin{quote}
\emph{He had an accomplice who was waiting for him in the car.}
\end{quote}

Also an adjective: \emph{un sourire complice} — a knowing smile.

\begin{cousinweb}
\emph{complice} is a learned borrowing from Late Latin \emph{complex}, \emph{complicem} \textquotedblleft{}partner, one folded in with\textquotedblright{} (EtymDB), built on \emph{plicāre} \textquotedblleft{}to fold\textquotedblright{}.
\end{cousinweb}

\begin{grammarlens}[title={What you've built}]
One more word for telling, explaining and describing at B1.
\end{grammarlens}

\subsection*{Your turn}
\begin{itemize}
\raggedright
  \item \emph{Say it:} \emph{le complice}
  \item \emph{Say it:} \emph{le complice}, once more
  \item \emph{Say it:} a sentence of your own with \emph{le complice}
  \item \emph{Recall:} say the French for to confess, then the French for the confession, then say \emph{le complice} again
\end{itemize}

\begin{tcolorbox}[breakable,colback=teal!4,colframe=teal!35!black,title={Before you move on}]
What does le complice mean? (\textquotedblleft{}The accomplice\textquotedblright{}.) Say it once more.
\end{tcolorbox}
\section[acquitter]{acquitter — to acquit; to pay off}

\label{lesson:FR-C295-acquitter}



\begin{quote}
Before the new one: say the French for the confession, then the French for the accomplice.
\end{quote}

\subsection*{You'll want to know: acquitter}
\textbf{acquitter} — \textquotedblleft{}to acquit; to pay off\textquotedblright{}.

\begin{quote}
Le tribunal a acquitté le jeune homme.
\end{quote}

\begin{quote}
\emph{The court acquitted the young man.}
\end{quote}

Also to pay a debt: \emph{s'acquitter d'une dette}.

\begin{cousinweb}
\emph{acquitter} is built on \emph{quitte} \textquotedblleft{}free of debt\textquotedblright{}, which comes from Medieval Latin \emph{quitus}, from Latin \emph{quiētus} \textquotedblleft{}at rest\textquotedblright{} (Lewis \& Short).
\end{cousinweb}

\begin{grammarlens}[title={What you've built}]
One more word for telling, explaining and describing at B1.
\end{grammarlens}

\subsection*{Your turn}
\begin{itemize}
\raggedright
  \item \emph{Say it:} \emph{acquitter}
  \item \emph{Say it:} \emph{acquitter}, once more
  \item \emph{Say it:} a sentence of your own with \emph{acquitter}
  \item \emph{Recall:} say the French for the confession, then the French for the accomplice, then say \emph{acquitter} again
\end{itemize}

\begin{tcolorbox}[breakable,colback=teal!4,colframe=teal!35!black,title={Before you move on}]
What does acquitter mean? (\textquotedblleft{}To acquit; to pay off\textquotedblright{}.) Say it once more.
\end{tcolorbox}
\section[relâcher]{relâcher — to release, to let go; to loosen}

\label{lesson:FR-C295-relacher}



\begin{quote}
Before the new one: say the French for the accomplice, then the French for to acquit.
\end{quote}

\subsection*{You'll want to know: relâcher}
\textbf{relâcher} — \textquotedblleft{}to release, to let go; to loosen\textquotedblright{}.

\begin{quote}
La police l'a relâché le lendemain.
\end{quote}

\begin{quote}
\emph{The police let him go the next day.}
\end{quote}

\emph{Relâcher ses efforts} — to slacken off.

\begin{cousinweb}
\emph{relâcher} is inherited through Old French \emph{relaschier} (EtymDB) from Latin \emph{relaxāre} \textquotedblleft{}to loosen\textquotedblright{} (Lewis \& Short).
\end{cousinweb}

\begin{grammarlens}[title={What you've built}]
One more word for telling, explaining and describing at B1.
\end{grammarlens}

\subsection*{Your turn}
\begin{itemize}
\raggedright
  \item \emph{Say it:} \emph{relâcher}
  \item \emph{Say it:} \emph{relâcher}, once more
  \item \emph{Say it:} a sentence of your own with \emph{relâcher}
  \item \emph{Recall:} say the French for the accomplice, then the French for to acquit, then say \emph{relâcher} again
\end{itemize}

\begin{tcolorbox}[breakable,colback=teal!4,colframe=teal!35!black,title={Before you move on}]
What does relâcher mean? (\textquotedblleft{}To release, to let go; to loosen\textquotedblright{}.) Say it once more.
\end{tcolorbox}
